\documentclass[11pt]{article}
\usepackage[a4paper,margin=25mm]{geometry}
\usepackage{fontspec,polyglossia}
\defaultfontfeatures{Extension=.otf}
\setmainfont{NimbusRoman}[UprightFont=*-Regular,BoldFont=*-Bold,ItalicFont=*-Italic,BoldItalicFont=*-BoldItalic]
\setsansfont{NimbusSans}[UprightFont=*-Regular,BoldFont=*-Bold,ItalicFont=*-Italic,BoldItalicFont=*-BoldItalic]
\setmonofont{NimbusMonoPS}[UprightFont=*-Regular,BoldFont=*-Bold,ItalicFont=*-Italic,BoldItalicFont=*-BoldItalic,Scale=0.85]
\newfontfamily\cyrillicfont{NimbusRoman}[UprightFont=*-Regular,BoldFont=*-Bold,ItalicFont=*-Italic,BoldItalicFont=*-BoldItalic,Script=Cyrillic]
\newfontfamily\cyrillicfontsf{NimbusSans}[UprightFont=*-Regular,BoldFont=*-Bold,ItalicFont=*-Italic,BoldItalicFont=*-BoldItalic,Script=Cyrillic]
\newfontfamily\cyrillicfonttt{NimbusMonoPS}[UprightFont=*-Regular,BoldFont=*-Bold,ItalicFont=*-Italic,BoldItalicFont=*-BoldItalic,Script=Cyrillic,Scale=0.85]
\def\ruCode{ru}
\ifx\wrehLang\ruCode
 \setdefaultlanguage{russian}\setotherlanguage{english}
 \newcommand{\Lang}[2]{#2}
\else
 \setdefaultlanguage{english}\setotherlanguage{russian}
 \newcommand{\Lang}[2]{#1}
\fi
\usepackage{amsmath,amssymb,amsthm,mathtools}
\usepackage{microtype,booktabs,tabularx,array,enumitem,xcolor}
\usepackage[hidelinks,unicode]{hyperref}
\newcolumntype{Y}{>{\raggedright\arraybackslash}X}
\newtheorem{definition}{\Lang{Definition}{Определение}}[section]
\newtheorem{proposition}[definition]{\Lang{Proposition}{Утверждение}}
\newtheorem{theorem}[definition]{\Lang{Theorem}{Теорема}}
\newtheorem{corollary}[definition]{\Lang{Corollary}{Следствие}}
\theoremstyle{definition}
\newtheorem{example}[definition]{\Lang{Example}{Пример}}
\newtheorem{remark}[definition]{\Lang{Remark}{Замечание}}
\newcommand{\Ccal}{\mathfrak C}
\newcommand{\Fcal}{\mathfrak F}
\newcommand{\Ical}{\mathfrak I}
\newcommand{\supp}{\operatorname{supp}}
\newcommand{\rank}{\operatorname{rank}}
\newcommand{\row}{\operatorname{row}}
\newcommand{\one}{\mathbf 1}
\definecolor{wrehblue}{RGB}{22,71,105}
\setlength{\emergencystretch}{2em}
\renewcommand{\arraystretch}{1.15}
\setlist{itemsep=3pt,topsep=4pt}
\hypersetup{pdftitle={\Lang{WR-V: Experiment and Realization Selection}{WR-V: Эксперимент и отбор реализаций}},pdfauthor={A. A. Malachevsky},pdfsubject={WREH; mathematical draft v0.1; unpeer-reviewed},pdfkeywords={experiment, instruments, transition identification, contextuality, nonnegative inverse problem}}
\title{\color{wrehblue}\textbf{WR-V}\\[6pt]
\Lang{Experiment and Realization Selection:\\Branch, Contextualize, Select,\\or Remain a Class}{Эксперимент и отбор реализаций:\\ветвление, контекстуализация, выбор\\или сохранение класса}\\[9pt]
\large\Lang{Mathematical draft v0.1}{Математический draft v0.1}}
\author{A. A. Malachevsky\\\small\href{https://orcid.org/0009-0008-6009-3196}{ORCID: 0009-0008-6009-3196}\\\small World Realizability \& Epistemic Horizons (WREH)}
\date{\Lang{8 October 2026}{8 октября 2026 года}}
